% Part: computability
% Chapter: recursive-functions
% Section: pr-functions

\documentclass[../../../include/open-logic-section]{subfiles}

\begin{document}

\olfileid{cmp}{rec}{prf}
\olsection{الدوال العودية البدائية}

فلنضبط الآن الوجهين اللذين نستخرج بهما دوال جديدة من دوال معلومة:
العودية البدائية والتركيب.

\begin{defn}
  \ollabel{defn:primitive-recursion} لتكن $\OLId{latin-ordinary}{f}$ دالة ذات $\OLId{latin-ordinary}{k}$ حجج
  ($\OLId{latin-ordinary}{k}\ge \OLMathNumeral{1}$)، ولتكن $\OLId{latin-ordinary}{g}$ ذات $\OLId{latin-ordinary}{k}+\OLMathNumeral{2}$ حجة. فالمعرّفة
  \emph{بالعودية البدائية من $\OLId{latin-ordinary}{f}$ و$\OLId{latin-ordinary}{g}$} هي الدالة~$\OLId{latin-ordinary}{h}$ ذات $\OLId{latin-ordinary}{k}+\OLMathNumeral{1}$ حجة،
  المعينة بالمعادلتين
  \begin{align*}
    \OLId{latin-ordinary}{h}(\OLId{latin-ordinary}{x}_\OLMathNumeral{0},\dots,\OLId{latin-ordinary}{x}_{\OLId{latin-ordinary}{k}-\OLMathNumeral{1}},\OLMathNumeral{0}) & =  \OLId{latin-ordinary}{f}(\OLId{latin-ordinary}{x}_\OLMathNumeral{0},\dots,\OLId{latin-ordinary}{x}_{\OLId{latin-ordinary}{k}-\OLMathNumeral{1}}) \\
    \OLId{latin-ordinary}{h}(\OLId{latin-ordinary}{x}_\OLMathNumeral{0},\dots,\OLId{latin-ordinary}{x}_{\OLId{latin-ordinary}{k}-\OLMathNumeral{1}},\OLId{latin-ordinary}{y}+\OLMathNumeral{1}) & = \OLId{latin-ordinary}{g}(\OLId{latin-ordinary}{x}_\OLMathNumeral{0},\dots,\OLId{latin-ordinary}{x}_{\OLId{latin-ordinary}{k}-\OLMathNumeral{1}}, \OLId{latin-ordinary}{y}, \OLId{latin-ordinary}{h}(\OLId{latin-ordinary}{x}_\OLMathNumeral{0},\dots,\OLId{latin-ordinary}{x}_{\OLId{latin-ordinary}{k}-\OLMathNumeral{1}}, \OLId{latin-ordinary}{y}))
  \end{align*}
\end{defn}

\begin{defn}
  \ollabel{defn:composition} لتكن $\OLId{latin-ordinary}{f}$ ذات $\OLId{latin-ordinary}{k}$ حجج، ولنأخذ
  $\OLId{latin-ordinary}{g}_\OLMathNumeral{0}$، …،~$\OLId{latin-ordinary}{g}_{\OLId{latin-ordinary}{k}-\OLMathNumeral{1}}$، وهي $\OLId{latin-ordinary}{k}$ دوال لكل منها $\OLId{latin-ordinary}{n}$ حجة. فالدالة الحاصلة \emph{بتركيب $\OLId{latin-ordinary}{f}$ مع $\OLId{latin-ordinary}{g}_\OLMathNumeral{0}$، …،~$\OLId{latin-ordinary}{g}_{\OLId{latin-ordinary}{k}-\OLMathNumeral{1}}$} هي الدالة
  ~ $\OLId{latin-ordinary}{h}$ ذات $\OLId{latin-ordinary}{n}$ حجة التي تعرفها
  \[
  \OLId{latin-ordinary}{h}(\OLId{latin-ordinary}{x}_\OLMathNumeral{0}, \dots, \OLId{latin-ordinary}{x}_{\OLId{latin-ordinary}{n}-\OLMathNumeral{1}}) =
  \OLId{latin-ordinary}{f}(\OLId{latin-ordinary}{g}_\OLMathNumeral{0}(\OLId{latin-ordinary}{x}_\OLMathNumeral{0}, \dots, \OLId{latin-ordinary}{x}_{\OLId{latin-ordinary}{n}-\OLMathNumeral{1}}), \dots, \OLId{latin-ordinary}{g}_{\OLId{latin-ordinary}{k}-\OLMathNumeral{1}}(\OLId{latin-ordinary}{x}_\OLMathNumeral{0}, \dots, \OLId{latin-ordinary}{x}_{\OLId{latin-ordinary}{n}-\OLMathNumeral{1}})).
  \]
\end{defn}

ونجعل في عداد الدوال العودية البدائية، مع $\Succ$ ودوال الإسقاط
\[
\Proj{\OLId{latin-ordinary}{n}}{\OLId{latin-ordinary}{i}}(\OLId{latin-ordinary}{x}_\OLMathNumeral{0},\dots,\OLId{latin-ordinary}{x}_{\OLId{latin-ordinary}{n}-\OLMathNumeral{1}}) = \OLId{latin-ordinary}{x}_\OLId{latin-ordinary}{i},
\]
لكل عدد طبيعي $\OLId{latin-ordinary}{n}$ ولكل $\OLId{latin-ordinary}{i}<\OLId{latin-ordinary}{n}$، الدالة~$\Zero(\OLId{latin-ordinary}{x})=\OLMathNumeral{0}$ أيضًا.

\begin{defn}
  الدوال العودية البدائية دوال من $\Nat^\OLId{latin-ordinary}{n}$ إلى
  $\Nat$، وتُضبط مجموعتها بالتعريف الاستقرائي الآتي:
  \begin{enumerate}
  \item $\Zero$ عودية بدائية.
  \item $\Succ$ عودية بدائية.
  \item كل دالة إسقاط $\Proj{\OLId{latin-ordinary}{n}}{\OLId{latin-ordinary}{i}}$ عودية بدائية.
  \item إذا كانت $\OLId{latin-ordinary}{f}$ دالة عودية بدائية ذات $\OLId{latin-ordinary}{k}$ حجج وكانت $\OLId{latin-ordinary}{g}_\OLMathNumeral{0}$،
    …،~$\OLId{latin-ordinary}{g}_{\OLId{latin-ordinary}{k}-\OLMathNumeral{1}}$ دوال عودية بدائية ذات $\OLId{latin-ordinary}{n}$ حجة، فإن تركيب $\OLId{latin-ordinary}{f}$ مع
    $\OLId{latin-ordinary}{g}_\OLMathNumeral{0}$، …،~$\OLId{latin-ordinary}{g}_{\OLId{latin-ordinary}{k}-\OLMathNumeral{1}}$ عودي بدائي.
  \item إذا كانت $\OLId{latin-ordinary}{f}$ دالة عودية بدائية ذات $\OLId{latin-ordinary}{k}$ حجج وكانت $\OLId{latin-ordinary}{g}$ دالة
    عودية بدائية ذات $\OLId{latin-ordinary}{k}+\OLMathNumeral{2}$ حجة، فإن الدالة المعرفة بالعودية البدائية
    من $\OLId{latin-ordinary}{f}$ و$\OLId{latin-ordinary}{g}$ عودية بدائية.
\end{enumerate}
\end{defn}

\begin{explain}
وحاصل التعريف أن مجموعة الدوال العودية البدائية أصغر مجموعة تضم
$\Zero$ و$\Succ$ ودوال الإسقاط~$\Proj{\OLId{latin-ordinary}{n}}{\OLId{latin-ordinary}{j}}$، وتكون منغلقة تحت
التركيب والعودية البدائية.

ولنا أن نصف هذه المجموعة على وجه آخر، بأن نبنيها
«مرحلة» بعد مرحلة. فلتكن $\OLId{latin-looped}{S}_\OLMathNumeral{0}$ مجموعة الدوال التي نبدأ بها، وهي $\Zero$ و$\Succ$ ودوال
الإسقاط. فهذه دوال المرحلة~$\OLMathNumeral{0}$. فإذا فرغنا من
تعريف $\OLId{latin-looped}{S}_\OLId{latin-ordinary}{i}$، جعلنا $\OLId{latin-looped}{S}_{\OLId{latin-ordinary}{i}+\OLMathNumeral{1}}$ مجموعة ما يحصل بإجراء
تركيب واحد أو عودية بدائية واحدة على دوال من~$\OLId{latin-looped}{S}_\OLId{latin-ordinary}{i}$.
فيكون
\[
\OLId{latin-looped}{S} = \bigcup_{\OLId{latin-ordinary}{i} \in \Nat} \OLId{latin-looped}{S}_\OLId{latin-ordinary}{i}
\]
مجموعة جميع الدوال العودية البدائية.
\end{explain}

ولنطبّق هذا الضبط فنثبت أن $\Add$ دالة عودية بدائية.

\begin{prop}
  دالة الجمع $\Add(\OLId{latin-ordinary}{x},\OLId{latin-ordinary}{y})=\OLId{latin-ordinary}{x}+\OLId{latin-ordinary}{y}$ عودية بدائية.
\end{prop}

\begin{proof}
قد سبق تعريف $\Add$ بالعودية البدائية من دالتين $\OLId{latin-ordinary}{f}$ و~$\OLId{latin-ordinary}{g}$،
وجاء التعريف على صورة \olref{defn:primitive-recursion}:
\begin{align*}
  \Add(\OLId{latin-ordinary}{x}_\OLMathNumeral{0}, \OLMathNumeral{0}) & = \OLId{latin-ordinary}{f}(\OLId{latin-ordinary}{x}_\OLMathNumeral{0}) = \OLId{latin-ordinary}{x}_\OLMathNumeral{0}\\
  \Add(\OLId{latin-ordinary}{x}_\OLMathNumeral{0}, \OLId{latin-ordinary}{y}+\OLMathNumeral{1}) & = \OLId{latin-ordinary}{g}(\OLId{latin-ordinary}{x}_\OLMathNumeral{0}, \OLId{latin-ordinary}{y}, \Add(\OLId{latin-ordinary}{x}_\OLMathNumeral{0}, \OLId{latin-ordinary}{y})) =
  \Succ(\Add(\OLId{latin-ordinary}{x}_\OLMathNumeral{0}, \OLId{latin-ordinary}{y}))
\end{align*}
فلا يبقى لإثبات عودية $\Add$ البدائية إلا إثباتها لـ$\OLId{latin-ordinary}{f}$ و$\OLId{latin-ordinary}{g}$. أما الأولى فإن
$\OLId{latin-ordinary}{f}(\OLId{latin-ordinary}{x}_\OLMathNumeral{0})=\OLId{latin-ordinary}{x}_\OLMathNumeral{0}=\Proj{\OLMathNumeral{1}}{\OLMathNumeral{0}}(\OLId{latin-ordinary}{x}_\OLMathNumeral{0})$، ودوال الإسقاط تعد عودية بدائية، فتكون
$\OLId{latin-ordinary}{f}$ عودية بدائية. أما $\OLId{latin-ordinary}{g}$ فهي الدالة الثلاثية الحجج $\OLId{latin-ordinary}{g}(\OLId{latin-ordinary}{x}_\OLMathNumeral{0},\OLId{latin-ordinary}{y},\OLId{latin-ordinary}{z})$
المعرفة بـ
\[
\OLId{latin-ordinary}{g}(\OLId{latin-ordinary}{x}_\OLMathNumeral{0}, \OLId{latin-ordinary}{y}, \OLId{latin-ordinary}{z}) = \Succ(\OLId{latin-ordinary}{z}).
\]
ولا يكفي هذا وحده للحكم بأن $\OLId{latin-ordinary}{g}$ عودية بدائية؛ فإن $\OLId{latin-ordinary}{g}$ غير $\Succ$
من جهة عدد الحجج: فـ$\Succ$ أحادية، و$\OLId{latin-ordinary}{g}$ لا بد أن تكون ثلاثية
الحجج. وإنما يتم تعريف $\OLId{latin-ordinary}{g}$ على الوجه الرسمي بالتركيب الآتي:
\[
\OLId{latin-ordinary}{g}(\OLId{latin-ordinary}{x}_\OLMathNumeral{0}, \OLId{latin-ordinary}{y}, \OLId{latin-ordinary}{z}) = \Succ(\Proj{\OLMathNumeral{3}}{\OLMathNumeral{2}}(\OLId{latin-ordinary}{x}_\OLMathNumeral{0}, \OLId{latin-ordinary}{y}, \OLId{latin-ordinary}{z})).
\]
وحيث إن $\Succ$ و$\Proj{\OLMathNumeral{3}}{\OLMathNumeral{2}}$ عوديتان بدائيتان، فإن $\OLId{latin-ordinary}{g}$
عودية بدائية أيضًا، لكونها تركيبًا منهما.
\end{proof}

\begin{prop}
  \ollabel{prop:mult-pr}
  دالة الضرب $\Mult(\OLId{latin-ordinary}{x},\OLId{latin-ordinary}{y})=\OLId{latin-ordinary}{x}\cdot \OLId{latin-ordinary}{y}$ عودية بدائية.
\end{prop}

\begin{proof}
  تمرين.
\end{proof}

\begin{prob}
  برهن \olref[cmp][rec][prf]{prop:mult-pr}: ضع تعريف $\Mult$
  بالعودية البدائية في الصورة المشترطة في
  \olref[cmp][rec][prf]{defn:primitive-recursion}، ثم بيّن أن الدالتين
  اللتين تقومان فيه مقام $\OLId{latin-ordinary}{f}$ و~$\OLId{latin-ordinary}{g}$ عوديتان بدائيتان.
\end{prob}

\begin{ex}
ولنعد إلى أول ما أوردناه مثالًا للتعريف بالعودية البدائية:
\begin{align*}
\OLId{latin-ordinary}{h}(\OLMathNumeral{0}) & =  \OLMathNumeral{1} \\
\OLId{latin-ordinary}{h}(\OLId{latin-ordinary}{y}+\OLMathNumeral{1}) & =  \OLMathNumeral{2} \cdot \OLId{latin-ordinary}{h}(\OLId{latin-ordinary}{y}).
\end{align*}
لا يندرج هذا التعريف كما هو في الصورة المشترطة في
\olref{defn:primitive-recursion}، لأن $\OLId{latin-ordinary}{k}=\OLMathNumeral{0}$. ويتضمن التعريف أيضًا
الثابتين $\OLMathNumeral{1}$ و$\OLMathNumeral{2}$. فأما الإشكال الأول فندفعه بإدخال حجة صورية، فنعرّف
الدالة~$\OLId{latin-ordinary}{h}'$ هكذا:
\begin{align*}
\OLId{latin-ordinary}{h}'(\OLId{latin-ordinary}{x}_\OLMathNumeral{0}, \OLMathNumeral{0}) & =  \OLId{latin-ordinary}{f}(\OLId{latin-ordinary}{x}_\OLMathNumeral{0}) = \OLMathNumeral{1} \\
\OLId{latin-ordinary}{h}'(\OLId{latin-ordinary}{x}_\OLMathNumeral{0}, \OLId{latin-ordinary}{y}+\OLMathNumeral{1}) & =  \OLId{latin-ordinary}{g}(\OLId{latin-ordinary}{x}_\OLMathNumeral{0}, \OLId{latin-ordinary}{y}, \OLId{latin-ordinary}{h}'(\OLId{latin-ordinary}{x}_\OLMathNumeral{0}, \OLId{latin-ordinary}{y})) = \OLMathNumeral{2} \cdot \OLId{latin-ordinary}{h}'(\OLId{latin-ordinary}{x}_\OLMathNumeral{0}, \OLId{latin-ordinary}{y}).
\end{align*}
يمكن تعريف الدالة $\OLId{latin-ordinary}{f}(\OLId{latin-ordinary}{x}_\OLMathNumeral{0})=\OLMathNumeral{1}$ من $\Succ$ و$\Zero$ بالتركيب:
$\OLId{latin-ordinary}{f}(\OLId{latin-ordinary}{x}_\OLMathNumeral{0})=\Succ(\Zero(\OLId{latin-ordinary}{x}_\OLMathNumeral{0}))$. ويمكن تعريف الدالة $\OLId{latin-ordinary}{g}$ بالتركيب من
$\OLId{latin-ordinary}{g}'(\OLId{latin-ordinary}{z})=\OLMathNumeral{2}\cdot \OLId{latin-ordinary}{z}$ ومن دوال الإسقاط:
\begin{align*}
\OLId{latin-ordinary}{g}(\OLId{latin-ordinary}{x}_\OLMathNumeral{0}, \OLId{latin-ordinary}{y}, \OLId{latin-ordinary}{z}) & = \OLId{latin-ordinary}{g}'(\Proj{\OLMathNumeral{3}}{\OLMathNumeral{2}}(\OLId{latin-ordinary}{x}_\OLMathNumeral{0}, \OLId{latin-ordinary}{y}, \OLId{latin-ordinary}{z}))
\intertext{ويمكن بدورها تعريف $g'$ بالتركيب هكذا:}
\OLId{latin-ordinary}{g}'(\OLId{latin-ordinary}{z}) & = \Mult(\OLId{latin-ordinary}{g}''(\OLId{latin-ordinary}{z}), \Proj{\OLMathNumeral{1}}{\OLMathNumeral{0}}(\OLId{latin-ordinary}{z}))
\intertext{حيث}
\OLId{latin-ordinary}{g}''(\OLId{latin-ordinary}{z}) & = \Succ(\OLId{latin-ordinary}{f}(\OLId{latin-ordinary}{z})),
\end{align*}
و$\OLId{latin-ordinary}{f}$ كما سبق: $\OLId{latin-ordinary}{f}(\OLId{latin-ordinary}{z})=\Succ(\Zero(\OLId{latin-ordinary}{z}))$. وحيث حصلت لنا~$\OLId{latin-ordinary}{h}'$،
فلنستعمل التركيب مرة أخرى ونضع
$\OLId{latin-ordinary}{h}(\OLId{latin-ordinary}{y})=\OLId{latin-ordinary}{h}'(\Proj{\OLMathNumeral{1}}{\OLMathNumeral{0}}(\OLId{latin-ordinary}{y}),\Proj{\OLMathNumeral{1}}{\OLMathNumeral{0}}(\OLId{latin-ordinary}{y}))$. فقد استخرجنا $\OLId{latin-ordinary}{h}$
من الدوال الأساسية بتركيبات وعوديات بدائية متتابعة؛ وبذلك ثبت أن
$\OLId{latin-ordinary}{h}$ عودية بدائية.
\end{ex}

\end{document}
